\chapter{Discussion}
\label{chap:discussion}

\section{Two regimes, two answers}
\label{sec:discussion-regimes}
Our central empirical pattern is a contrast between regimes. In-domain on
ASL Citizen~\cite{aslcitizen} --- a large, signer-independent benchmark
(crowdsourced, 52 signers) whose training split already exposes the model
to many signers, and whose participant demographics have since been
documented and released for bias-mitigation
research~\cite{atwellStudyingMitigatingBiases2024} --- appearance
interventions bought
little: +1.1 R@1 on a frozen-feature head, nothing at the backbone. In the
low-data NGT setting, where the model must bridge a real signer gap from a
handful of recordings, motion-preserving synthetic counterfactuals produced
a clear gain. The simplest reading is that appearance invariance is not a
property a model has or lacks in the abstract; it is relative to the
diversity already present in training. The Logos checkpoint we
use~\cite{ovodovLogosWellTemperedPretrain2025} is trained jointly on three
isolated-sign corpora (Logos, AUTSL, and WLASL); the Logos dataset alone
contributes roughly 200{,}000 videos from 305 signers --- and the other two corpora add 36 and 119 further signers on top, respectively. A
backbone with that exposure has already
consumed the invariance that our in-domain augmentations could have
provided; the NGT model, seeing four signers, had not.
% Signer count verified against the Logos paper §3.1 (2026-07-11):
% 199,668 videos, 381 signers, 2,863 lemmas; "the most extensive available
% ISLR dataset by the number of signers.
%The paper reports 381 signers total and constructs a signer-disjoint 80/20 split, selecting about 20% of signer identities exclusively for testing. That gives:
%approximately 305 signers in its training set and 76 in its test set."

A second factor compounds this signer-diversity gap for the raw, untrained
backbone specifically. That Logos transfers well onto ASL Citizen even
before any ASL-Citizen-specific training (\S\ref{sec:discussion-signclip})\todome{this ref
points at the SignCLIP discussion section, which is currently an empty
todo stub --- either fill that section or re-anchor this claim to the
retrieval table directly}
makes sense once the pretraining recipe is considered: Logos is trained
jointly with WLASL, and WLASL and ASL Citizen share a large fraction of
their vocabulary --- both are American Sign Language isolated-sign corpora
drawn from an overlapping sign inventory, so the backbone has effectively
already seen much of what ASL Citizen asks of it. The shift to NGT removes
that overlap on three independent axes at once: NGT is a different sign
language entirely (Dutch Sign Language rather than American), its
recordings come from a different capture setup (signers in motion-capture
suits in a studio, rather than webcam video), and its unit of analysis is
sentence-level continuous signing rather than the isolated, gloss-level
clips both ASL Citizen and WLASL provide. Any one of these differences
would weaken transfer; together they are a plausible reason why the same
pretrained backbone that carries over to ASL Citizen with no
target-corpus training needs
the full synthetic-appearance intervention to bridge the NGT signer gap.

This reading also explains why our ASL Citizen null differs from SDDA's positive
result~\cite{fuSignerDiversitydrivenData2024b} on Phoenix2014T: their
diffusion-augmented contrastive setup operates on a studio-recorded corpus
with far less signer variety than ASL Citizen --- PHOENIX-2014T contains nine signers in
total~\cite{fuSignerDiversitydrivenData2024b} ---  i.e.\ closer to our NGT
% Nine signers verified 2026-07-12 in three independent sources: Camgoz et
% al. 2018 §4 (primary; no bib entry yet — optional to add), the SDDA paper
% itself ("includes data from only nine different signers"), and the Isharah
% paper's comparison table. Cited to SDDA since it is already in the bib and
% is the paper whose corpus is being discussed.
regime than to our ASL Citizen regime.
Notably, the SDDA authors themselves
flag the low visual realism of their synthesized signers as a limitation;
our Flux9B pipeline produces markedly more realistic edits,
yet realism alone did not rescue the in-domain intervention --- evidence
that the binding constraint is the regime, not image quality.
% Piotr's question (2026-07-12): "could we directly compare to SDDA by
% training the model on Phoenix?" Answer, per the resolution report: not as
% a small job — SDDA is a sign language TRANSLATION result (BLEU/ROUGE on
% continuous PHOENIX-2014T), so a like-for-like comparison needs the full
% continuous-recognition/SLT stack this thesis scopes out; our Flux pipeline
% could generate the augmentations, but the training/eval harness does not
% exist in this project. Left out of the draft text; revisit as Future Work
% if desired.

SDDA~\cite{fuSignerDiversitydrivenData2024b} is evaluated exclusively on
PHOENIX-2014T, and not on that corpus's standard split: following the signer-independent protocol introduced for
this corpus by Jin and Zhao, the data are re-partitioned into four settings
in which one of signers 3, 4, 7 and 8 is held out as the test set, the three
data-poor signers 2, 6 and 9 form the validation set, and the remaining five
signers supply the training data (their Appendix~B, Table~7).
% Verified 2026-07-29 in the SDDA paper: §4.1 "the data of signers 3, 4, 7,
% and 8 constitute the test set, while the remaining data serve as the
% training or validation set. Given the relatively limited data for signers
% 2, 6, and 9, we allocate their corresponding data to the validation set";
% Table 7 lists train signers {1,4,5,7,8}/{1,3,5,7,8}/{1,3,4,5,8}/{1,3,4,5,7}
% = five per split, with 7163/6639/6980/6880 samples, dev 411, test
% 683/1207/866/966. Split totals equal 8257 = the full corpus.
What separates their setting from ours is the diversity available during training, and five real
signers places SDDA far closer to our NGT regime than to ASL
Citizen.

Recording conditions are held fixed --- one studio, one camera
setup, one weather-forecast register --- so the shift their augmentation has
to bridge is signer appearance and signing style alone, narrower than the
corpus and viewpoint changes our NGT evaluation layers on top of it. How
much the held-out-signer choice matters on this corpus is documented
independently: under a full signer-disjoint re-evaluation of
PHOENIX-2014T, translation scores roughly halve and vary by a factor of
${\sim}4$ across the choice of held-out
signer~\cite{artiagaRethinkingSignLanguage2025}.\todome{reference verified 04.08: Artiaga et al.\ (Findings of EMNLP
2025) reports signer-disjoint BLEU-4 drops 21.44$\to$3.59 (GFSLT-VLP),
15.74$\to$8.26 (GASLT), 22.74$\to$3.66 (SignCL) --- supports the
sentence; bib metadata matches the ACL Anthology record}
% Cross-check (Artiaga et al., "Rethinking Sign Language Translation: The
% Impact of Signer Dependence on Model Evaluation", Findings of EMNLP 2025,
% aclanthology.org/2025.findings-emnlp.997): confirms the PHOENIX-2014T
% DEFAULT split is signer-overlapping ("Its default split includes
% overlapping signers across ... nine signers", §1) --- which SDDA avoids.
% Their §3.1 re-splits into nine folds (one test signer, one dev signer,
% seven train signers; train 5,100--7,893 videos). Table 1: default-split
% BLEU-4 21.44/15.74/22.74 (GFSLT-VLP/GASLT/SignCL) falls to a 9-fold mean
% of 10.53/10.24/4.18; per-fold GFSLT-VLP BLEU-4 spans 3.59--17.30, i.e.
% the choice of held-out signer moves the score by ~4x. Table 2: 1.0--3.3%
% of test sentences also occur in train even under signer folds. Note their
% §3.1 prose mis-states the default sizes as 7,022/269/966 (those are their
% Fold 8); Table 1 gives the correct 7,096/519/642. No bib entry yet.


\section{Does SignCLIP-style text alignment help or hurt?}
\label{sec:discussion-signclip}
% RECOMMENDATION (Claude 04.08): WRITE, do not remove. Removal is not a
% local operation: RQ1 names SignCLIP, the WLASL diagnosis probes it, the
% whole in-domain head experiment IS a SignCLIP head, the head signer-probe
% panel, six of nine t-SNE panels, and the appendix analogy/gender/iconicity
% analyses are computed in SignCLIP-head spaces; sec:discussion-regimes
% cross-references this section. Draft below; adopt/trim/reject.
\todome{DRAFT START (Claude 04.08): section 1 of 2. Strictly restates
what tab:asl\_citizen\_retrieval, tab:asl\_citizen\_t2v,
tab:head\_signer\_probes and fig:tsne\_grid already show; no new numbers,
no new citations. The open-vocabulary/anchor-neighbourhood argument cut
on 03.08 (see the todome at line 1905) is deliberately NOT
reconstructed here.}

For task quality, the answer in our high-diversity regime is: it helps,
modestly, and it is the only place the in-domain synthetic data moved
anything. A SignCLIP-style head trained from scratch on frozen Logos
features already reaches
0.750 % line 1829 (from-scratch head, V2T R@1)
V2T R@1 over 2{,}731 glosses % line 1854 (caption: 2731 unique glosses)
(random-baseline Median R $1366$), % line 1853-1854 (caption)
which is the concrete sense in which the pretrained Logos backbone,
itself never trained on ASL Citizen by its
authors~\cite{ovodovLogosWellTemperedPretrain2025}, already carries the
benchmark: the head is the only target-corpus learning involved
(\S\ref{sec:discussion-regimes} points here for exactly this claim).
On top of that, fine-tuning the head adds real-data gains
(0.789), % line 1831 (Fine-tuned)
adding the Flux edits as plain extra training items adds the
$+1.1$-point gain quoted in \S\ref{sec:discussion-regimes}
(0.800), % line 1833 (Fine-tuned + synth. data); +1.1 vs 0.789 at line 1831
and carrying the same recipe onto fine-tuned backbones tops out at
0.818. % line 1843 (FT+synth. recipe on native)
The text-to-video direction repeats the ordering, with the combination
strongest on every metric
(P@1 $0.919$, % line 1908 (Fine-tuned + synth. data, T2V)
$0.925$--$0.935$ on fine-tuned backbones), % lines 1912-1917 (recipe-on-FT rows)
and the head also tightens the gloss geometry that the backbone leaves
loose: on the phonologically confusable t-SNE families, every head sits
at silhouette $s{=}0.50$--$0.51$ against $0.31$ for the frozen backbone
alone. % lines 1989-1994 (fig:tsne_grid caption)
None of this is a large effect on an already-strong backbone --- the
gains are one to three R@1 points --- but the direction is consistently
positive, and nothing in the retrieval tables suggests that adding a
text tower costs accuracy.

What text alignment does to \emph{signer identity} is more interesting
than what it does to accuracy, because the two probe readouts pull
apart. Relative to the backbone spaces, every head makes signers
\emph{more} geometrically separated at fixed gloss
($\rho$ of $22.0$--$26.3\times$ % lines 2076-2080 (tab:head_signer_probes)
against $17.67\times$ for the frozen reference) % line 2075 (frozen row)
while making identity far \emph{less} linearly decodable
($20$--$25\%$ against $66.9\%$ on the same train split). % lines 2090-2092 (caption)
% (35 signers, chance 0.029 -- caption lines 2085-2087)
We resist reading a mechanism into this beyond what the probes state:
the head reorganises the space so that whatever signer information
survives is no longer linearly accessible, while the coarse prototype
geometry, if anything, stretches. Which of the two readouts one
privileges decides whether text alignment ``helped'' with signer
invariance, and \S\ref{sec:results-ft-probes} argues neither should be
read as a proxy for the other. What can be said flatly is that text
alignment did not remove signer structure --- no objective in this
thesis did --- it moved where that structure is measurable.

The clearest ``hurt'' is not the text alignment itself but what happens
when the counterfactuals are pushed into the head as explicit
positives: pair-SupCon at $K{=}3$ drops V2T R@1 to
0.567 % line 1849 (K=3 pair-SupCon)
against
0.800 % line 1833 (Fine-tuned + synth. data)
for the same edits used as plain additional data, and $K{=}5$ falls
further to
0.414. % line 1850 (K=5 pair-SupCon; note the checkpoint caveat todo at line 1873)
The same synthetic videos therefore help or hurt depending purely on
the objective through which they enter --- as data, a small gain; as
forced positives, a large loss --- mirroring the high-$K$ degradation
in the NGT ladder (\S\ref{sec:future-work} discusses the candidate
discriminative fix). Our overall answer is accordingly narrow: a
SignCLIP-style text-aligned head on frozen features is a cheap,
slightly-better-than-neutral addition in this regime, it is the only
component through which the in-domain synthetic data produced any gain,
and it changes the signer-identity story qualitatively rather than
erasing it.

\todome{DRAFT END: caveats. (i) The K=5 pair-SupCon number comes from
checkpoint\_last.pt with no best checkpoint saved --- there is an open
todo at line 1873 to remove the K=5 row; if it goes, drop the ``and
$K{=}5$ falls further to 0.414'' clause here (the K=3 contrast carries
the point alone). (ii) The three published pose-row numbers
(0.39/0.46/0.60, lines 1836-1838) are deliberately not used here
pending the verification todo at lines 1874-1876. (iii) The head-probe
comparison to 66.9\% is train-split, 35 signers --- do not mix it with
the 75.3\% test-split figure in prose. (iv) If Piotr prefers a shorter
section, paragraph 1 + paragraph 3 suffice; paragraph 2 (the
dissociation) overlaps most with \S\ref{sec:results-ft-probes} and can
be compressed to two sentences.}


\section{RGB versus pose, revisited}
\label{sec:discussion-rgb-pose}

\todome{DRAFT START (Claude 04.08): section 2 of 2. Reuses only
citation keys already present in the file (Intro pose-caveat block,
lines 191-240; Related Work input-modality section, lines 414-424) and
only numbers printed in sec:results-l2-raw and the WLASL probes.}

Pose input is widely adopted in SLR partly \emph{because} it is assumed
to discard appearance: the robustness of pose-based models is
attributed to inherent invariance to clothing, background and
lighting~\cite{moryossefEvaluatingImmediateApplicability2021,
kratimenosIndependentSignLanguage2020, tranRegionAwarePoseModeling2025},
% keys as used at lines 423-424
and the SignEval organisers went as far as distributing
\emph{only} pose features for a signer-independent task, citing signer
overfitting and privacy~\cite{luqmanSignEval2025Challenge2025a,
alyamiIsharahLargeScaleMultiScene2026b}. % keys as used at lines 182-183
Our diagnostic probes complicate the premise behind both moves. On ASL
Citizen, mean-pooled MediaPipe keypoints are the \emph{most}
signer-separable of the three raw feature spaces we test, on both
readouts: their inter/intra distance ratio is
$29.9\times$ [$28.4$, $32.0$] % table line 1360 (tab:raw_feature_l2_probe)
against
$22.7\times$ and $22.2\times$ % lines 1361-1362 (Logos, I3D)
for the two RGB backbones --- confidence intervals overlapping
neither's % prose lines 1372-1374
--- and a linear head recovers the signer at
$98.7\%$ % line 1360 (MediaPipe decoding; chance 3.0%, caption line 1354
% and prose lines 1370-1371)
against $94.2\%$ (I3D) and $65.9\%$ (Logos), an ordering unchanged on
the signer-disjoint test split
($99.4\%$/$95.2\%$/$74.3\%$, chance $9.1\%$). % lines 1374-1376
Camera framing cannot carry this: the keypoints are root-centred at the
shoulder midpoint and scale-normalised to shoulder width, so what
identifies the signer is shoulder-normalised geometry --- body
proportions, habitual posture, movement style. % prose lines 1376-1381
Nor is it gloss content in disguise: in the leave-one-gloss-out WLASL
probe, raw MediaPipe features still retrieve the correct signer at
$0.207$ top-1 (chance $0.013$), % line 1296 (tab:cross_gloss_signer_probe)
and the original pose-based SignCLIP model's own embedding space does so
\emph{more} strongly, at $0.302$. % line 1295
Pose discards pixels; on this evidence it retains --- and arguably
distills --- precisely the person-specific signal the modality is
assumed to remove.

This lands on the same side as the fairness literature the Introduction
reviews, from the opposite direction. That literature's caveat is about
the \emph{extractor}: pose pipelines are themselves learned RGB models
whose training and evaluation sets under-represent darker-skinned
individuals~\cite{lachanceCaseStudyFairness2023,
hasanHi5SyntheticData2024}, % keys as used at line 194
whose audits reach conclusions that flip with the evaluation
set~\cite{lachanceCaseStudyFairness2023,
xiangFairHumanCentricImage2025}, % keys as used at line 203
whose hand tracking degrades exactly where signing is linguistically
dense~\cite{moryossefOptimizingHandRegion2024,
moryossefEvaluatingImmediateApplicability2021}, % keys as used at lines 218-221
and whose ASL Citizen pose baseline shows larger skin-tone and gender
disparities than its RGB
counterpart~\cite{atwellStudyingMitigatingBiases2024}. % key as used at line 210
Our probes add the complementary caveat about the \emph{representation}:
even when extraction succeeds and coordinates are normalised, the
resulting features are more signer-identifying than either RGB space we
compare against. Together the two caveats suggest that treating pose
distribution as an appearance-bias or privacy fix --- as in the
pose-only challenge protocol above --- deserves the same scrutiny as
any other debiasing claim: the appearance axis a pose skeleton removes
(clothing, background, skin pixels) is not the only axis along which
signers are identifiable, and by our measurements not the dominant one.

The scope of this claim should be stated as narrowly as the evidence.
Ours is a bias probe, not a matched-capacity recognition comparison: we
never train a pose-based recogniser ourselves (the pose rows of
Table~\ref{tab:asl_citizen_retrieval} are the published results of the
original SignCLIP model), % setup lines 1058-1061
so we make no claim about whether pose-based \emph{models} exploit this
identity signal, how much it costs them under signer shift, or how the
modalities compare at matched training --- questions the recent
RGB-versus-pose results
cited in \S\ref{sec:related-work-modality}~\cite{minCloserLookSkeletonbased2025,
wongSignRepEnhancingSelfSupervised2025} % keys as used at lines 419, 240
address from the accuracy side. The comparison is also asymmetric in
pretraining: the MediaPipe features are untrained coordinates while the
RGB features come from supervised backbones, and the Logos column shows
exactly what such supervision does --- gloss-supervised pretraining
sheds linearly accessible identity ($65.9\%$ against I3D's $94.2\%$)
while leaving the distance geometry untouched. % prose lines 1382-1390
What the probes do establish is the narrower, assumption-level point:
signer identity is not removed by pose extraction by construction, and
a pipeline that adopts pose \emph{for} invariance or privacy inherits
an unexamined assumption rather than a guarantee.

\todome{DRAFT END: caveats. (i) The WLASL cross-gloss probe prose
(lines 1315-1319) itself hedges that vocabulary bias may inflate those
numbers --- if that hedge matters to you here, add ``(with the
vocabulary-bias caveat of \S\ref{sec:results-diagnosis})'' after the
0.207/0.302 sentence. (ii) The ``not the dominant one'' clause at the
end of paragraph 2 rests on the both-probes ordering in
tab:raw\_feature\_l2\_probe; if that feels too strong, weaken to ``and on
our probes not even the larger one''. (iv) This section overlaps the Intro's pose-caveat paragraph (lines
191-240) by design --- Discussion may restate Intro framing --- but if
Piotr wants zero repetition, the second paragraph can be cut to its
first and last sentences.}
